% MURM: A Local-First, Provider-Agnostic Swarm Intelligence Engine
% for Calibrated Societal Prediction
%
% Template for: arXiv preprint / NeurIPS Workshop on Social Simulation /
%               ICML Workshop on Foundation Models for Decision Making
%
% To compile: pdflatex paper.tex && bibtex paper && pdflatex paper.tex x2

\documentclass[10pt,twocolumn]{article}
\usepackage[utf8]{inputenc}
\usepackage[T1]{fontenc}
\usepackage{times}
\usepackage{geometry}
\geometry{left=1.5cm,right=1.5cm,top=2cm,bottom=2cm}
\usepackage{hyperref}
\usepackage{amsmath}
\usepackage{booktabs}
\usepackage{graphicx}
\usepackage{listings}
\usepackage{xcolor}
\usepackage{natbib}

\lstset{
  basicstyle=\small\ttfamily,
  breaklines=true,
  frame=single,
  backgroundcolor=\color{gray!8}
}

\title{MURM: A Local-First, Provider-Agnostic Swarm Intelligence Engine\\
for Calibrated Societal Prediction}

\author{
  [Author Names]\\
  [Institution]\\
  \texttt{[contact email]}
}

\date{}

\begin{document}
\maketitle

\begin{abstract}
Multi-agent LLM simulation has emerged as a promising approach to societal prediction, yet existing systems exhibit three systematic shortcomings: dependency on proprietary external services for knowledge graph storage, absence of agent diversity enforcement leading to measurable herd behavior, and no statistical framework for calibrating or quantifying uncertainty in predictions. We present \textsc{MURM}, an open-source swarm intelligence prediction engine that addresses all three gaps. Our contributions are: (1) a fully local knowledge graph pipeline using NetworkX and ChromaDB that eliminates external service dependencies; (2) a quota-sampled, seeded agent population generator that enforces specified opinion distributions across five-level Likert scales, demonstrably reducing premature consensus relative to unconstrained LLM generation; (3) per-round emergence metrics including Shannon opinion entropy, Gini posting inequality, and a polarization index derived from entropy decay; (4) a multi-seed sensitivity analysis framework that produces distributional predictions with calibration infrastructure based on Brier scoring. We evaluate \textsc{MURM} on [N] prediction tasks spanning [domains] and show that [key result]. All code, data, and evaluation scripts are released at \url{https://github.com/[org]/murm}.
\end{abstract}

\section{Introduction}

The hypothesis underlying swarm intelligence prediction systems is that simulating the micro-level interactions of a diverse agent population produces macro-level emergent dynamics that track real-world collective behavior more faithfully than individual-agent or statistical baselines \cite{park2023generative, gao2023s3, mou2024agentsocially}. Several recent systems instantiate this hypothesis \cite{mirofish2025, oasis2024}, but they share a cluster of engineering and methodological weaknesses that limit their scientific utility.

First, all existing open systems we are aware of rely on Zep Cloud or equivalent proprietary services for knowledge graph storage. This creates a reproducibility barrier: the graph state is opaque, quota-limited, and not inspectable. Second, agent populations are generated by prompting an LLM to produce $N$ diverse personas in a single call — a procedure that produces clusters rather than distributions. The OASIS paper \cite{oasis2024} documents that LLM agents polarize significantly faster than empirical social media users, a finding consistent with the homogeneity of unconstrained LLM-generated populations. Third, existing systems produce point-estimate predictions with no confidence bounds, no cross-run variance analysis, and no mechanism for longitudinal calibration against ground truth.

\textsc{MURM} addresses each of these in turn. We make the following concrete contributions:

\begin{itemize}
  \item A two-pass entity extraction pipeline (ontology generation followed by constrained extraction) grounded in a locally-stored, fully-inspectable NetworkX knowledge graph.
  \item A quota-sampled persona generator that enforces a target opinion distribution (normal, bimodal, power-law, or uniform) using the largest-remainder method, guaranteeing exact distribution targets independent of LLM tendencies.
  \item A per-round emergence metrics collector measuring Shannon entropy, Gini coefficient, opinion velocity, and polarization index — the first systematic attempt to quantify emergent dynamics in LLM social simulation.
  \item A sensitivity analysis framework that runs $k$ independent seeds and produces a distributional prediction with variance labeling and a Brier score infrastructure for downstream calibration.
\end{itemize}

\section{Related Work}

\paragraph{LLM Agent Simulation.} \citet{park2023generative} introduced Generative Agents, demonstrating coherent long-horizon behavior in a sandbox environment. \citet{gao2023s3} extended this to social network simulation. OASIS \cite{oasis2024} scales to $10^6$ agents for Twitter dynamics simulation, at the cost of significant engineering complexity and documented herd behavior. MiroFish \cite{mirofish2025} wraps OASIS in a product interface but inherits its methodological weaknesses.

\paragraph{Opinion Dynamics.} The Deffuant model \cite{deffuant2000mixing} and Axelrod's culture model \cite{axelrod1997dissemination} established theoretical frameworks for opinion convergence and polarization. Recent work \cite{bail2018exposure, brady2021MAD} quantifies these dynamics empirically. Our entropy-based metrics are designed to detect the same phenomena in LLM simulations.

\paragraph{Calibration in Predictive Systems.} Forecasting systems such as Metaculus and PredictionBook maintain calibration curves mapping stated confidence to empirical accuracy. \citet{guo2017calibration} showed that neural networks are systematically miscalibrated. We apply Brier scoring \cite{brier1950verification} as a foundational calibration metric for swarm predictions.

\section{System Design}

\subsection{Knowledge Graph Construction}

Given seed document $D$, we run a two-pass extraction:

\textbf{Pass 1 — Ontology generation.} We prompt the LLM to derive a minimal domain-specific schema $\mathcal{O} = (E, R)$ where $E$ is a set of entity type labels and $R$ is a set of relation type labels, constrained to types evidenced in $D$. Temperature is set to 0.1 to minimize hallucination.

\textbf{Pass 2 — Entity extraction.} We prompt the LLM to extract all entities and relations from $D$ using $\mathcal{O}$ as a hard constraint. Entities not matching any $e \in E$ are rejected. Relations referencing non-extracted entities are filtered by a post-processing step.

The resulting graph $G = (V, \mathcal{E})$ is persisted locally as a NetworkX \texttt{DiGraph} with JSON serialization. Semantic search is provided via ChromaDB with sentence-transformer embeddings (\texttt{all-MiniLM-L6-v2}).

\subsection{Diversity-Enforced Population Generation}

Let $N$ be the target population size and $\mathbf{p} = (p_1, \ldots, p_5)$ be the target opinion distribution over the five-level scale (strongly agree, agree, neutral, disagree, strongly disagree). We compute integer quotas using the largest-remainder method:

\begin{equation}
  c_i = \left\lfloor N p_i \right\rfloor + \mathbb{1}\left[i \in \text{top-}r \text{ remainders}\right]
\end{equation}

where $r = N - \sum_i \lfloor N p_i \rfloor$ is the residual count. This guarantees $\sum_i c_i = N$ exactly and eliminates the LLM's tendency to cluster around neutral opinions.

Each agent is generated individually with its assigned opinion and influence role visible in the prompt, preventing the LLM from overriding the assignment. Influence role weights follow an empirical social network topology: 5\% influencers, 25\% active users, 55\% passive users, 10\% skeptics, 5\% amplifiers \cite{bakshy2012role}.

\subsection{Simulation Engine}

The engine is fully asynchronous. Each round $t$:

\begin{enumerate}
  \item A subset $A_t \subset \mathcal{A}$ of agents is selected to act based on their \texttt{reaction\_speed} parameter sampled against a seeded RNG.
  \item Each agent in $A_t$ receives the environment context feed $F_t$ and up to $k=3$ semantically relevant graph entities retrieved by querying the ChromaDB embedder with the agent's last action summary.
  \item The agent produces an action (post, reply, share, or abstain) and optionally updates its opinion position.
  \item The action is appended to the environment feed and written to the JSONL trace.
\end{enumerate}

Counterfactual events can be injected at any specified round, inserting content into the environment feed as if from an external source. This enables the "what if" analysis described in MiroFish's README but not implemented there.

\subsection{Emergence Metrics}

At each round $t$ we compute:

\textbf{Opinion entropy.} Let $\mathbf{q}^{(t)} = (q_1^{(t)}, \ldots, q_5^{(t)})$ be the empirical opinion distribution at round $t$. Shannon entropy:
\begin{equation}
  H(t) = -\sum_{i=1}^5 q_i^{(t)} \log_2 q_i^{(t)}
\end{equation}
Maximum entropy is $\log_2 5 \approx 2.32$ bits (uniform distribution). $H(t) \to 0$ indicates full consensus.

\textbf{Polarization index.}
\begin{equation}
  P(t) = \frac{\log_2 5 - H(t)}{\log_2 5} \in [0, 1]
\end{equation}

\textbf{Gini coefficient.} Let $a_i^{(t)}$ be the number of actions taken by agent $i$ at round $t$. Gini:
\begin{equation}
  G(t) = \frac{2 \sum_{i=1}^{|A_t|} i \cdot a_{(i)}^{(t)}}{|A_t| \sum_i a_i^{(t)}} - \frac{|A_t|+1}{|A_t|}
\end{equation}
where $a_{(i)}^{(t)}$ is the $i$-th order statistic of action counts.

\textbf{Opinion velocity.}
\begin{equation}
  V(t) = \frac{1}{|\mathcal{A}|} \sum_{a \in \mathcal{A}} |o_a^{(t)} - o_a^{(t-1)}|
\end{equation}
where $o_a^{(t)} \in \{-2, -1, 0, 1, 2\}$ is agent $a$'s opinion value at round $t$.

\subsection{Calibration and Sensitivity Analysis}

For a prediction question $Q$, the system generates a point prediction $\hat{y}$ with stated confidence $c \in [0,1]$.

\textbf{Multi-seed sensitivity.} We run $k$ independent simulations with seeds $s_1, \ldots, s_k$. For outcome variable $X$ (e.g., final entropy), we compute $\bar{X}$ and $\sigma_X$. Variance is labeled as low ($\sigma_X < 0.05$), medium ($0.05 \leq \sigma_X < 0.15$), or high ($\sigma_X \geq 0.15$).

\textbf{Brier score.} When ground truth $y \in \{0,1\}$ is observed:
\begin{equation}
  BS = (c - y)^2
\end{equation}
$BS = 0$ is perfect; the no-skill baseline for binary events is $BS = 0.25$.

\section{Experiments}

\subsection{Experimental Setup}

[Describe your experimental setup: prediction domains, number of runs, models used, comparison baselines.]

\subsection{Herd Behavior Reduction}

We compare the opinion entropy at convergence between populations generated with and without diversity enforcement.

\begin{table}[h]
  \centering
  \caption{Final opinion entropy (bits) by population generation method and distribution target. Higher entropy indicates less premature convergence.}
  \begin{tabular}{lcc}
    \toprule
    Method & Mean $H(T)$ & Std $H(T)$ \\
    \midrule
    Unconstrained (MiroFish-style) & [value] & [value] \\
    Quota-sampled Normal & [value] & [value] \\
    Quota-sampled Bimodal & [value] & [value] \\
    \bottomrule
  \end{tabular}
  \label{tab:entropy}
\end{table}

\subsection{Prediction Accuracy}

[Results table: Brier scores across domains, compared to baseline LLM-only predictions.]

\subsection{Emergence Dynamics}

[Time series plots of $H(t)$, $P(t)$, $V(t)$ for representative runs. Describe what patterns are visible.]

\subsection{Counterfactual Analysis}

[Show that injecting a counterfactual event at round $t^*$ produces statistically detectable changes in $H(t)$ for $t > t^*$.]

\section{Limitations}

\textbf{LLM-as-Human validity.} LLM agents do not faithfully replicate human cognitive processes. Opinion shifts in our simulations reflect LLM decision boundaries, not human psychology. Results should be interpreted as reasoning over stated priors, not as behavioral models.

\textbf{Context window budget.} Each agent's action prompt is bounded by the LLM's context window. Long simulations accumulate more environment content than can fit; our system uses a fixed-size recent feed which may lose long-range dependencies.

\textbf{Calibration requires longitudinal data.} Brier scores can only be computed once real-world outcomes resolve. We provide the infrastructure but cannot present calibration curves without a historical prediction dataset, which is future work.

\textbf{Computational cost.} $N$ agents across $T$ rounds with $k$ seeds requires $O(NTk)$ LLM calls. At $N=50$, $T=30$, $k=3$, this is 4,500 calls per prediction question. Cost estimates are surfaced to users pre-run.

\section{Conclusion}

We presented \textsc{MURM}, a swarm intelligence prediction engine that addresses the three principal methodological weaknesses of existing systems: external service dependency, unconstrained agent generation, and absence of calibration infrastructure. The system is fully local, provider-agnostic, and ships with quantitative emergence metrics and a sensitivity analysis framework suitable for research use. We release the code, paper, and evaluation data to support reproducibility and further work on calibrated collective intelligence.

\bibliographystyle{plainnat}
\bibliography{references}

\end{document}
